\documentclass[12pt,a4paper]{article}
\usepackage[a4paper,left=25mm,right=20mm,top=24mm,bottom=24mm]{geometry}
\usepackage{fontspec}
\usepackage{polyglossia}
\setdefaultlanguage{russian}
\setotherlanguage{english}
\IfFontExistsTF{Libertinus Serif}{
    \setmainfont{Libertinus Serif}
    \newfontfamily\cyrillicfont{Libertinus Serif}
}{
    \IfFontExistsTF{DejaVu Serif}{
        \setmainfont{DejaVu Serif}
        \newfontfamily\cyrillicfont{DejaVu Serif}
    }{
        \setmainfont{Times New Roman}
        \newfontfamily\cyrillicfont{Times New Roman}
    }
}
\IfFontExistsTF{DejaVu Sans Mono}{\setmonofont{DejaVu Sans Mono}[Scale=0.82]}{}
\usepackage{microtype}
\usepackage{setspace}
\setstretch{1.20}
\usepackage{parskip}
\setlength{\parindent}{0pt}
\usepackage{booktabs}
\usepackage{tabularx}
\usepackage{longtable}
\usepackage{array}
\usepackage{ragged2e}
\newcolumntype{P}[1]{>{\RaggedRight\arraybackslash}p{#1}}
\newcolumntype{Y}{>{\RaggedRight\arraybackslash}X}
\setlength{\emergencystretch}{3em}
\usepackage{enumitem}
\usepackage{needspace}
\usepackage{amsmath}
\usepackage{amssymb}
\setlist{topsep=3pt,itemsep=2pt,parsep=0pt}
\usepackage{graphicx}
\usepackage{xcolor}
\definecolor{codegray}{RGB}{62,66,72}
\usepackage{listings}
\lstset{
    basicstyle=\ttfamily\footnotesize,
    breaklines=true,
    breakatwhitespace=false,
    columns=fullflexible,
    showstringspaces=false,
    frame=single,
    rulecolor=\color{black!20},
    xleftmargin=4pt,
    xrightmargin=4pt,
    aboveskip=8pt,
    belowskip=8pt
}
\usepackage{tikz}
\usetikzlibrary{arrows.meta,positioning,shapes.geometric}
\usepackage{hyperref}
\hypersetup{hidelinks,pdfauthor={Agibalov Vladislav},pdflang={ru-RU}}
\usepackage{fancyhdr}
\pagestyle{fancy}
\setlength{\headheight}{14pt}
\fancyhf{}
\fancyhead[L]{\small mangen}
\fancyhead[R]{\small Техническая документация · 2026}
\fancyfoot[C]{\thepage}
\renewcommand{\headrulewidth}{0.2pt}
\setcounter{tocdepth}{2}
\newcommand{\impl}{\textbf{Реализация.}~}

\title{mangen\\Программа и результаты испытаний}
\author{Агибалов Владислав Сергеевич\\НИЯУ МИФИ, ИВТ}
\date{2026 год}

\begin{document}
\maketitle
\begin{abstract}
Документ устанавливает воспроизводимую методику проверки C17-приложения
\texttt{mangen}, включая unit/CTest, системное вычисление SHA-2, испытания
доступа к файлам, многопоточность, sanitizers, утечки и переносимость.
Отдельно указаны результаты, сообщённые в журнале запуска macOS,
и проверки, не подтверждённые на Linux для полного проекта.
\end{abstract}
\tableofcontents
\newpage

\section{Цель и критерии приёмки}

Приёмка требует, чтобы выход был однозначен, сортирован и полностью
совпадал с независимыми SHA-2, многопоточность не меняла его содержимого,
ошибки доступа не скрывались, все ресурсы освобождались и сборка выполнялась
на заявленных UNIX-платформах без скачивания зависимостей.

\begin{longtable}{P{0.14\textwidth}P{0.47\textwidth}P{0.29\textwidth}}
\toprule
ID & Метод & Условие успеха \\
\midrule
\endhead
V01 & CTest со всеми библиотеками & Все запущенные тесты пройдены \\
V02 & \texttt{BUILD\_TESTING=OFF} & Сборка только C, без GoogleTest \\
V03 & \texttt{cmake --install} & Бинарник в \texttt{bin/mangen} \\
V04 & Полное доступное дерево & Exit 0, пустой stderr \\
V05 & Дерево с chmod 000 & Exit 1, полный доступный набор \\
V06 & Системные SHA-256/384/512 & Побайтовое совпадение строк \\
V07 & \texttt{-j 1} и \texttt{-j 4} & Идентичный stdout \\
V08 & ASan/UBSan/LSan & Нет выявленных нарушений \\
V09 & Компиляторы и ОС & Проходит матрица на macOS/Linux \\
V10 & Подпись и контроль ZIP & Изменение байта отвергается \\
\bottomrule
\end{longtable}

\section{Протоколы выполненных испытаний}

В разделе приведены результаты запуска программы на macOS и результаты
отдельных проверок сборки на Linux. Числовые итоги на macOS получены
из предоставленного журнала команд и их вывода; для Linux объём
фактически выполненных операций указан отдельно.

\Needspace{8\baselineskip}
\subsection{Результаты macOS по журналу}

\begin{longtable}{P{0.48\textwidth}P{0.43\textwidth}}
\toprule
Измерение & Наблюдение \\
\midrule
\endhead
CTest, обычная сборка & 163 из 163 тестов пройдены \\
CTest, ASan + UBSan & 163 из 163 тестов пройдены \\
Большое дерево с ограничениями & Exit 1; 590 строк stdout; 248 строк диагностики \\
Системная проверка SHA-256 & 590 из 590 строк совпали \\
После восстановления разрешений & 6401 запись; exit 0; диагностик нет \\
Повторная системная проверка & 6401 из 6401 строк SHA-256 совпали \\
Сравнение потоков & stdout с одним и четырьмя потоками идентичен \\
Запуск ASan/UBSan на полном дереве & Exit 0; 6401 записей; stderr пуст \\
Apple \texttt{leaks} & Сводка: 0 утечек, 0 leaked bytes; есть предупреждение об ограничениях доступа к процессу \\
Наблюдавшийся пиковый RSS & Около 3.45 MiB в приведённом прогоне \\
\bottomrule
\end{longtable}

Пределы интерпретации: успешный CTest подтверждает конкретную сборку,
компилятор, опции и платформу, а не математическое отсутствие всех гонок;
\texttt{leaks} сообщил об ограниченном режиме инспекции, поэтому «0 утечек»
не является полным доказательством отсутствия всех утечек во всех ветвях.

\section{Структура контролируемых деревьев}

\texttt{scripts/generate\_test\_trees.sh} создаёт в отдельном системном
временном каталоге два дерева. Каждое включает 6400 небольших текстовых
файлов с различным содержимым, один бинарный файл 16 MiB, один пустой файл,
одну ссылку на регулярный файл и один пустой каталог. Итог для доступного
дерева --- 6403 записей. Во втором дереве для файлов и директорий случайно
назначаются режимы chmod; не менее одного файла и каталога принудительно
остаются недоступными владельцу.

\begin{lstlisting}
TEST_ROOT=$(bash scripts/generate_test_trees.sh)
bash scripts/verify_test_trees.sh "$TEST_ROOT" \
    ./build/install/bin/mangen
\end{lstlisting}

Полная проверка включает SHA-256, SHA-384, SHA-512 и режимы 1/4 потока
для обоих деревьев, то есть 12 запусков приложения. Выход каждого сравнивается
с эталоном на базе системных утилит; проверка не полагается только на
количество строк или синтаксис HEX.

\subsection{Права доступа}

Проверку необходимо запускать от обычного пользователя. В режиме root
системные привилегии могут обходить ограничения chmod. Соответствующие
утверждения теста в таком окружении пропускаются с явным предупреждением.
Успешный частичный запуск должен одновременно удовлетворять трём условиям:

\begin{enumerate}
\item код возврата равен 1;
\item stdout содержит в точности все доступные ожидаемые файлы;
\item stderr содержит сообщения об отказах доступа.
\end{enumerate}

Независимая системная \texttt{find} также может получить ошибку обхода;
это ожидаемо для защищённого дерева, но недопустимо для полного.

\section{Методика проверки переносимости}

\begin{lstlisting}
bash scripts/test_portability.sh
\end{lstlisting}

Скрипт определяет текущую ОС, ищет доступные C-компиляторы, отдельно
конфигурирует \texttt{BUILD\_TESTING=OFF} и \texttt{ON}, компилирует,
выполняет CTest, устанавливает приложение и проверяет его запуск.
Отсутствие зависимости C++ в режиме \texttt{OFF} является самостоятельным
критерием.

\begin{center}
\setlength{\tabcolsep}{4pt}
\begin{tabularx}{\textwidth}{P{0.32\textwidth}YYY}
\toprule
Платформа & Основная сборка & CTest & Sanitizers \\
\midrule
macOS & По журналу: да & 163/163 & ASan/UBSan: да \\
Linux x86\_64, библиотеки & C17 синтаксис: да & Не запускался & Не запускались \\
Linux, полный проект & Не подтверждено & Не подтверждено & Не подтверждено \\
\bottomrule
\end{tabularx}
\end{center}

Дополнительная проверка на Ubuntu 24.04.3 x86\_64 с GCC 13.3.0:
конфигурация Concurrency Kit 0.7.2 завершилась успешно; проверка
\texttt{-fsyntax-only} для доступных модулей \texttt{matcher},
\texttt{queue}, \texttt{catalog}, \texttt{hash}, \texttt{hash\_pool}
завершилась с кодом 0. Обнаружено одно предупреждение GCC о сравнении,
которое всегда ложно, в \texttt{queue.c}. Это не линковка и не запуск CTest:
полный CMake-цикл не выполнен (CMake отсутствовал в проверочной среде).

Системы macOS и Linux отличаются системными утилитами SHA, некоторыми
заголовками и feature-test macros, реализациями \texttt{pthread}, правилами
линковки и архитектурными компонентами Concurrency Kit. Окончательное
утверждение о Linux-поддержке требует исполнения именно в Linux, а не
переноса результатов macOS по аналогии.

\section{Диагностика памяти}

\begin{lstlisting}
bash scripts/test_memory.sh "$TEST_ROOT/clean"
\end{lstlisting}

Скрипт конфигурирует Debug-сборку с ASan и UBSan,
запускает CTest и приложение на полном дереве.
На Linux через \path{ASAN_OPTIONS} включается LeakSanitizer;
на macOS при наличии можно дополнительно запускать штатный \texttt{leaks}.

Ни один sanitizer не гарантирует нахождение всех ошибок. Для очереди MPMC
и её жизненного цикла дополнительно рекомендован ThreadSanitizer там,
где он поддерживается данной архитектурой и набором компиляторов.

\section{Подпись и тест отрицательного сценария}

Создать отдельный ключ подписанта, подготовить trusted\_signers со
сверенным публичным ключом, подписать манифест и выполнить проверку.
Затем изменить один байт в \emph{копии} манифеста, не меняя подпись:
\texttt{verify\_manifest\_signature.sh} должен вернуть ненулевой код.
Изменение файла \texttt{trusted\_signers} или неподписанного манифеста
не рассматривается как допустимый способ «исправить» ошибку подписи.

При выпуске ZIP аналогичный отрицательный тест проводится по копии архива:
после изменения одного байта \texttt{scripts/verify\_release.sh} должен
завершиться с ошибкой. Проверочный TXT должен быть получен независимо
от недоверенного канала, иначе злоумышленник может заменить и ZIP, и TXT.

\section{Отчёт по завершении приёмки}

Протокол каждого запуска должен содержать дату, ОС и архитектуру,
версии CMake/компилятора, полный Git commit, команды, количество файлов,
завершение процесса, stdout/stderr, сводку sanitizer, hash ZIP и
тест отрицательного сценария. Неуспешные и недоступные проверки не должны
обозначаться пройденными.

\addcontentsline{toc}{section}{Список литературы}
\begingroup
\footnotesize
\setlength{\parskip}{0pt}
\begin{thebibliography}{9}
\bibitem{googletest} GoogleTest. CMake Quickstart. \url{https://google.github.io/googletest/quickstart-cmake.html}
\bibitem{sanitizers} LLVM Clang. AddressSanitizer. \url{https://clang.llvm.org/docs/AddressSanitizer.html}
\bibitem{lsan} LLVM Clang. LeakSanitizer. \url{https://clang.llvm.org/docs/LeakSanitizer.html}
\bibitem{cmake} Kitware. CTest documentation. \url{https://cmake.org/cmake/help/latest/manual/ctest.1.html}
\bibitem{sha} NIST. FIPS 180-4. \url{https://csrc.nist.gov/pubs/fips/180-4/upd1/final}
\end{thebibliography}
\endgroup
\end{document}
